\documentclass[11pt]{article}
\usepackage[a4paper,margin=1in]{geometry}
\usepackage[T1]{fontenc}
\usepackage{lmodern,microtype}
\usepackage{amsmath,amssymb,amsthm,mathtools}
\usepackage{booktabs,enumitem}
\usepackage{xcolor}
\usepackage[numbers,sort&compress]{natbib}
\usepackage[colorlinks=true,linkcolor=blue!55!black,citecolor=blue!55!black]{hyperref}

\newtheorem{theorem}{Theorem}[section]
\newtheorem{proposition}[theorem]{Proposition}
\newtheorem{corollary}[theorem]{Corollary}
\theoremstyle{remark}
\newtheorem{remark}[theorem]{Remark}

\title{Monodromy Integrality Barrier for Signed Moment Fits}
\author{Bin Wang}
\date{July 2026}

\begin{document}
\maketitle

\begin{abstract}
After idempotent polynomial selectors collapse to binary masks, one may try
arbitrary signed shell weights.  On the expanded RH-254 windows, minimum-norm
real weights fit the order-$2$--$12$ anchor tolerance at all 32 endpoints.
Every fit is fractional, however, and maximum weights range from $39.64$ to
$3.02\times10^{11}$.  We prove that the weighted moment germ equals
$\prod_j(1-\lambda_jz)^{w_j}$ near zero and has monodromy
$e^{2\pi i w_j}$ around $\lambda_j^{-1}$.  Integer combined exponents are
therefore necessary for a single-valued meromorphic determinant quotient.
The 32 continuous signed fits are finite compressions, not legal selectors.
\end{abstract}

\section{Weighted moment germs}

Let $\lambda_1,\dots,\lambda_J$ be distinct nonzero complex numbers and let
$w_j\in\mathbb C$.  Define the local logarithmic moment germ
\begin{equation}
 \mathcal F_w(z)=\exp\left[-\sum_{n\ge1}\frac{z^n}{n}
 \sum_{j=1}^Jw_j\lambda_j^n\right]
 \label{eq:germ}
\end{equation}
for $|z|<\min_j|\lambda_j|^{-1}$.

\begin{theorem}[Monodromy integrality]
\label{thm:monodromy}
On the initial disk,
\begin{equation}
 \mathcal F_w(z)=\prod_{j=1}^J(1-\lambda_jz)^{w_j}.
 \label{eq:product}
\end{equation}
Analytic continuation once around $z=\lambda_j^{-1}$ multiplies the germ by
$e^{2\pi i w_j}$.  Hence $\mathcal F_w$ is single-valued meromorphic near all
reciprocal roots if and only if every combined exponent at a repeated root is
an integer.
\end{theorem}

\begin{proof}
The power-series identity
$-\sum_{n\ge1}(\lambda z)^n/n=\log(1-\lambda z)$ proves
\eqref{eq:product} on the initial disk.  Continuing the logarithm once around
its zero changes it by $2\pi i$; exponentiation therefore produces the stated
factor.  Trivial monodromy is equivalent to $e^{2\pi iw_j}=1$, or
$w_j\in\mathbb Z$.  Integer powers are rational meromorphic factors, proving
the converse.
\end{proof}

\begin{remark}[Necessary, not sufficient]
Integer weights are necessary for a determinant ratio, but an operator
realization also requires the corresponding algebraic multiplicities and
invariant spaces.  Negative integers describe denominator or superdeterminant
contributions; they are not automatically present in the noisy operator.
\end{remark}

\section{Expanded-window signed fit}

At each endpoint let $V$ be the real order-$2$--$12$ shell-power matrix and
$d$ the difference between the Perron/parity-removed trace and deterministic
anchor.  We solve the unconstrained minimum-norm problem
\begin{equation}
 \min\{\|w\|_2:Vw\ \text{is the least-squares approximation to }d\}.
 \label{eq:lstsq}
\end{equation}
This is deliberately more permissive than the RH-255 box
\citep{WangRH255}.

\begin{center}
\small
\begin{tabular}{lr}
\toprule
quantity & value\\
\midrule
endpoints within local tolerance & 32/32\\
full row-rank systems & 28/32\\
weighted residual range & $2.15\times10^{-15}$--0.003783\\
fractional weights per endpoint & 20--34\\
maximum absolute weight range & 39.64--$3.02\times10^{11}$\\
integer-weight fits & 0/32\\
maximum monodromy defect & 1.9999999996\\
\bottomrule
\end{tabular}
\end{center}

The four rank-deficient systems still meet their comparatively larger local
tolerances.  The smallest singular value across the batch is
$2.71\times10^{-19}$, which helps explain the coefficient explosion.  The
fit is therefore both nonintegral and numerically ill-conditioned.

\begin{corollary}[Scoped illegality of the continuous fits]
None of the 32 least-norm signed fits defines a single-valued finite
determinant quotient with the fitted shell exponents.
\end{corollary}

\begin{proof}
Every fitted vector contains at least 20 weights farther than $10^{-8}$ from
the integer lattice.  Apply Theorem~\ref{thm:monodromy}.
\end{proof}

\section{Boundary and next route}

This result does not exclude a bounded integer signed selector or a
non-product quotient-kernel cancellation.  It does exclude the inference
``finite signed fit implies determinant selector.''  RH-258 should audit a
small integer lattice and then separately ask for an operator realization.
Gates A--E remain false/open.  No Hilbert--Polya operator, Riemann-zero
identification, zeta-divisor equality, or RH implication is claimed.

\bibliographystyle{plainnat}
\bibliography{references}
\end{document}
